\documentclass[11pt]{article}

\usepackage[margin=1in]{geometry}
\usepackage{amsmath,amssymb}
\usepackage{hyperref}
\usepackage{listings}
\usepackage{xcolor}

\lstset{
  basicstyle=\ttfamily\footnotesize,
  keywordstyle=\color{blue},
  commentstyle=\color{gray},
  stringstyle=\color{orange},
  frame=single,
  language=Python
}

\title{Scientific Python}
\date{September 17, 2024}

\begin{document}

\maketitle
\tableofcontents

\section*{Introduction}

Python is an interpreted language. This makes writing a program easier. However, a consequence is that loops are slow. In scientific computing, it is common to work with large containers and to iterate over them. This requires a specific usage that is different from ``general Python'', and is made possible thanks to two packages:
\begin{itemize}
\item The \texttt{numpy} package,
\item The \texttt{pandas} package.
\end{itemize}

The \texttt{numpy} library provides essential containers in scientific computing, such as vectors or matrices, and supplies a whole set of mathematical functions that make it possible to avoid using loops. The \texttt{pandas} package, on the other hand, offers a specific container: the data table, and all the functions to manipulate and study it.

\section{Generalities}

\subsection{Installation}

Python is an interpreted language. It requires downloading the interpreter (what is generally called Python), which can be downloaded from the official website:  
\url{https://www.python.org/downloads/}.

However, it is not very practical to use it for development as is. It allows executing Python files directly and writing in the console, but this is hardly convenient. For a more user-friendly development, it is useful to use an integrated development environment (IDE). One may for example use RCode:  
\url{https://rcode.dev/}.

\subsection{First steps}

Python has weak typing: one does not specify the type, but all variables have one. This type is implicit.

Assignment of a variable is done with the operator \texttt{=}.

\begin{lstlisting}
x = 3
x
\end{lstlisting}

The function \texttt{type} makes it possible to obtain the type of a variable.

\begin{lstlisting}
x = "Python"
type(x)
<class 'str'>
type(True)
<class 'bool'>
type(1)
<class 'int'>
type(1.)
<class 'float'>
\end{lstlisting}

\subsection{Conditionals}

The type \texttt{bool} represents Booleans: the variable takes the value \texttt{True} or \texttt{False}.

The operator \texttt{and} corresponds to logical AND.

\begin{lstlisting}
True and False
\end{lstlisting}

The operator \texttt{or} corresponds to logical OR.

\begin{lstlisting}
True or False
\end{lstlisting}

The operator \texttt{not} corresponds to logical NOT.

\begin{lstlisting}
not True
\end{lstlisting}

A condition is written with \texttt{if} in the following way:

\begin{lstlisting}
if (condition):
    instructions
\end{lstlisting}

or without parentheses:

\begin{lstlisting}
if condition:
    instructions
\end{lstlisting}

The indentation above is fundamental. It defines the scope of the \texttt{if}. Indentation is part of the syntax: as long as the next line is indented, it belongs to the \texttt{if}.

The condition is a Boolean that takes the value \texttt{True} or \texttt{False}.

\begin{lstlisting}
x = 7  # This is a comment
if x % 2 == 1:  # If x modulo 2 equals 1
    print("x is odd")  # print writes in the console
\end{lstlisting}

It is possible to add an alternative using the keyword \texttt{else} followed by a colon.

\begin{lstlisting}
x = x + 1  # x = 8
if x % 2 == 1:  # FALSE
    print("x is odd")
else:
    print("x is even")
\end{lstlisting}

Once again, indentation defines the scope of the \texttt{else}.

\subsection{Loops}

A \texttt{for} loop is an instruction that is carried out over a set $i\in I$.

The semantics of the \texttt{for} loop is:
\begin{lstlisting}
for i in I:
    instructions
\end{lstlisting}

The most common case is to perform an instruction over
\[
i\in\{0,1,\dots,n-1\}
\]
for $n\in\mathbb{N}^*$.  
The function \texttt{range} evaluated at $n$ makes it possible to represent this set. The variable $I$ can therefore be \texttt{range(n)}.

\begin{lstlisting}
for i in range(4):
    print("Iteration i = " + str(i))  # str(i) converts i to a string
                                     # + concatenates strings
\end{lstlisting}

It is possible to start the iteration at some $n_0>0$ with \texttt{range(n0, n)}.

\medskip
The \texttt{while} loop makes it possible to repeat an instruction as long as a condition is satisfied. Its syntax is:

\begin{lstlisting}
while (condition):
    instructions
\end{lstlisting}

or without parentheses:

\begin{lstlisting}
while condition:
    instructions
\end{lstlisting}

A \texttt{for} loop can always be written as a \texttt{while} loop:
\begin{lstlisting}
i = 1
while i <= n:
    instructions
    i = i + 1
\end{lstlisting}

Be careful with infinite loops when using \texttt{while}.

\subsection{Functions}

A function takes one or several arguments (possibly none) and returns a variable (possibly none).

The syntax to define a function is:

\begin{lstlisting}
def function_name(arg1, ..., argn):
    instructions
\end{lstlisting}

The call is made with:

\begin{lstlisting}
function_name(arg1, ..., argn)
\end{lstlisting}

If the function returns a value, it ends with \texttt{return}; as soon as a \texttt{return} is encountered, the execution of the function stops.

Below is an example with the factorial function.

\begin{lstlisting}
def factorial(n):
    if n == 0:
        return 1
    else:
        f = 1
        for i in range(2, n+1):
            f *= i
        return f

factorial(0), factorial(1), factorial(5)
\end{lstlisting}

Above, \texttt{*=} multiplies the variable on the left of the operator by the element on the right while modifying it; the corresponding line is equivalent to \texttt{f = f*i}.

\subsection{Containers}

\subsubsection{Lists}

A list is a collection of elements without constraints. It can mix simple variables of any type with objects.

\begin{lstlisting}
l = [1, 4, "RCode"]
print(l)
\end{lstlisting}

Access to elements is done with \texttt{[i]} where $i$ ranges from $0$ to the size of the list minus one.

\begin{lstlisting}
[l[0], l[2]]
\end{lstlisting}

It is possible to include lists inside lists.

\begin{lstlisting}
h = [l, 8]
print(h)
\end{lstlisting}

Calling the first element of index $0$ returns the list \texttt{l}.

\begin{lstlisting}
h[0]
\end{lstlisting}

\subsubsection{Tuples}

Tuples are exactly like lists except that they are not modifiable.

\begin{lstlisting}
t = (1, 4, "RCode")
print(t)
\end{lstlisting}

One can even omit the parentheses:
\begin{lstlisting}
t = 1, 4, "RCode"
\end{lstlisting}

\subsubsection{Dictionaries}

Dictionaries are lists in which each member is named.

\begin{lstlisting}
c1 = "id"
d = {c1: 1, "x": 4, "label": "RCode"}
print(d)
\end{lstlisting}

Access is by name:
\begin{lstlisting}
d["label"]
\end{lstlisting}

\section{The numpy package}

In scientific computing, it is frequent to have arrays whose elements are all of the same type. The Python list is too flexible.

The \texttt{numpy} package introduces a dedicated array object called \texttt{numpy.ndarray} (vectors, matrices, etc.).

Always import numpy:

\begin{lstlisting}
import numpy as np
\end{lstlisting}

\subsection{The vector}

A vector is a collection of values of the same type.

We can create a vector with \texttt{np.array}:

\begin{lstlisting}
x = np.array([1, 3, 7])
x
\end{lstlisting}

To initialize an array of size $n$:

\begin{lstlisting}
n = 10
z = np.zeros(n)
z
\end{lstlisting}

Also \texttt{np.ones}, \texttt{np.empty}.

\texttt{np.arange(a,b)} creates integers on $[a,b[$:

\begin{lstlisting}
np.arange(0, 5)
\end{lstlisting}

A vector has an attribute \texttt{.size}:

\begin{lstlisting}
x.size
\end{lstlisting}

Indexing:

\begin{lstlisting}
x[0]
\end{lstlisting}

Boolean mask selection:

\begin{lstlisting}
x = np.array([1, 3, 5])
x[np.array([False, True, True])]
\end{lstlisting}

Index selection:

\begin{lstlisting}
x[np.array([0, 2, 2, 1])]
\end{lstlisting}

Slicing:

\begin{lstlisting}
x[0:2]
\end{lstlisting}

Deleting:

\begin{lstlisting}
np.delete(x, 1)
\end{lstlisting}

\medskip
\textbf{First fundamental rule.}  
All operators applied to two vectors of the same size return a vector where the operator has been applied element by element.

\begin{lstlisting}
y = np.arange(0, 3)
x + y
x * y
x ** y
x == y
\end{lstlisting}

\medskip
\textbf{Second fundamental rule.}  
All operators applied to a vector and a scalar return a vector where the operator has been applied to each element.

\begin{lstlisting}
2*x + 1
x ** 2
x % 2
x <= 2
\end{lstlisting}

\subsection{Aggregation functions}

An aggregation function takes a vector and returns a real number.

\medskip
Sum:
\[
\texttt{np.sum(x)} = \sum_{i=1}^{n} x_i .
\]

Mean:
\[
\texttt{np.mean(x)} = \frac{1}{n} \sum_{i=1}^{n} x_i .
\]

Variance:
\[
\texttt{np.var(x)} = \sum_{i=1}^{n} \left(x_i - \texttt{np.mean}(x)\right)^2 .
\]

Standard deviation:
\[
\texttt{np.std(x)} = \sqrt{\texttt{np.var(x)}} .
\]

\subsection{Vector language and efficient programming}

We say that a function $f$ is \emph{vectorial} if it takes one or several vectors as arguments (and possibly other arguments of size 1) and returns a vector where the function has been applied element by element.

\begin{lstlisting}
for i in range(x.size):
    x[i] = f(i, x[i], z[i])
\end{lstlisting}

If $f$ is vectorial:
\begin{lstlisting}
x = f(np.arange(x.size), x, z)
\end{lstlisting}

\medskip
\textbf{Remark.} Prefer \texttt{np.arange(x.size)} over \texttt{range(x.size)} when using \texttt{numpy}.

Example:
\begin{lstlisting}
np.sum(np.arange(n+1)**2)
\end{lstlisting}

Norm of $x+y$:
\begin{lstlisting}
np.sqrt(np.sum(x**2 + y**2))
\end{lstlisting}

Conditional update pattern:
\begin{lstlisting}
ind = h(np.arange(x.size), x, z)
x[ind] = f(np.arange(x.size)[ind], x[ind], z[ind])
\end{lstlisting}

Example for $y=(x-1)_+$:
\begin{lstlisting}
y = np.zeros(x.size)
ind = x >= 1
y[ind] = x[ind] - 1
\end{lstlisting}

\subsection{Simulation of random variables}

The \texttt{numpy} package uses a sequence of pseudo-random numbers generated by an arithmetic recurrence.

\begin{itemize}
\item The default algorithm used is the Mersenne--Twister. Its period is $2^{19937}-1$, which is approximately $10^{6000}$.
\item This algorithm is often considered to be the best compromise between efficiency and quality.
\item The generator is initialized automatically.
\item It is possible to choose a seed with \texttt{np.random.seed}.
\item This allows one to reproduce the same simulations for debugging purposes.
\end{itemize}

Uniform on $[0,1]$:
\begin{lstlisting}
np.random.random(size=None)
\end{lstlisting}

Example:
\begin{lstlisting}
n = 5
U = np.random.random(n)
U
\end{lstlisting}

Uniform on $[a,b]$:
\begin{lstlisting}
np.random.uniform(low=0., high=1., size=None)
\end{lstlisting}

Standard normal:
\begin{lstlisting}
np.random.randn(size=None)
\end{lstlisting}

Example:
\begin{lstlisting}
n = 5
Z = np.random.randn(n)
Z
\end{lstlisting}

General normal:
\begin{lstlisting}
np.random.normal(loc=0., scale=1., size=None)
\end{lstlisting}

Some distributions:
\begin{itemize}
\item Binomial $B(m,p)$: \texttt{np.random.binomial(m, p, n)}
\item Poisson $P(\lambda)$: \texttt{np.random.poisson(a, n)}
\item Geometric starting from $0$, $G(p)$: \texttt{np.random.geometric(p, n)}
\item Negative binomial $NB(m,p)$: \texttt{np.random.negative\_binomial(m, p, n)}
\item Log-normal $LN(m,s^2)$: \texttt{np.random.lognormal(m, s, n)}
\item Exponential $E(b)$: \texttt{np.random.exponential(1/b, n)}
\item Gamma $G(a,b)$: \texttt{np.random.gamma(a, 1/b, n)}
\end{itemize}

For densities/cdf/quantiles use \texttt{scipy.stats}:
\begin{lstlisting}
import scipy.stats as stats
x = np.linspace(-3., 3., 601)

stats.norm.pdf(x)
stats.norm.cdf(x)
a = np.linspace(0., 1., 21)
stats.norm.ppf(a)
\end{lstlisting}

\newpage

\subsection{Drawing graphs with matplotlib}

A simple and standard way to draw graphs in Python is to use the \texttt{matplotlib} package.
The usual import is:

\begin{lstlisting}
import matplotlib.pyplot as plt
\end{lstlisting}

\medskip
To plot a curve from data points \texttt{x} and \texttt{y} (same size):

\begin{lstlisting}
x = np.linspace(0., 2*np.pi, 200)
y = np.sin(x)

plt.figure()
plt.plot(x, y)
plt.xlabel("x")
plt.ylabel("sin(x)")
plt.title("A first plot")
plt.grid(True)
plt.show()
\end{lstlisting}

\medskip
It is possible to draw several curves on the same figure, and add a legend:

\begin{lstlisting}
y2 = np.cos(x)

plt.figure()
plt.plot(x, y, label="sin")
plt.plot(x, y2, label="cos")
plt.legend()
plt.grid(True)
plt.show()
\end{lstlisting}

\medskip
\textbf{Remark.} When studying performance, a log-scale is often useful:

\begin{lstlisting}
plt.figure()
plt.plot(n_values, times_vec, label="vectorized")
plt.plot(n_values, times_loop, label="for-loop")
plt.xscale("log")
plt.yscale("log")
plt.xlabel("n")
plt.ylabel("time (s)")
plt.legend()
plt.grid(True)
plt.show()
\end{lstlisting}

To save the figure to a file, use \texttt{plt.savefig("figure.png")} before \texttt{plt.show()}.

\subsection{Measuring computation time with the time module}

\paragraph{Remark (timing functions).}
To measure computation time in Python, one may use either \texttt{time.time()} or \texttt{time.perf\_counter()}.
The function \texttt{time.time()} returns the current \emph{wall-clock} time (a Unix timestamp), so it is convenient and often sufficient when the computation is long enough (e.g.\ a few tenths of a second or more).
However, it is not designed for precise benchmarking: its resolution can be limited, and the measurement may be affected by external factors such as system clock adjustments.
In contrast, \texttt{time.perf\_counter()} is a high-resolution \emph{monotonic} timer dedicated to performance measurements, and it typically provides more stable results when the code runs in only a few milliseconds (which is common for vectorized NumPy operations).
If you choose to use \texttt{time.time()}, it is recommended to repeat the measurement a few times and keep the minimum (or the median) to reduce measurement noise.

\begin{lstlisting}
import time

t0 = time.perf_counter()
# code to time
t1 = time.perf_counter()
print("Elapsed time:", t1 - t0, "seconds")
\end{lstlisting}

\medskip
A convenient pattern is to write a small helper function that measures the runtime of a function call:

\begin{lstlisting}
def time_call(func, *args, repeat=3):
    """Return a robust estimate of the runtime of func(*args)."""
    # warm-up (helps reduce first-call overhead)
    func(*args)

    tmin = float("inf")
    for _ in range(repeat):
        t0 = time.perf_counter()
        func(*args)
        t1 = time.perf_counter()
        tmin = min(tmin, t1 - t0)
    return tmin
\end{lstlisting}

\medskip
Finally, to compare a vectorized implementation with a loop implementation, one can measure times for increasing sizes and plot the curves:

\begin{lstlisting}
import numpy as np
import matplotlib.pyplot as plt

n_values = np.array([10**3, 3*10**3, 10**4, 3*10**4, 10**5, 3*10**5, 10**6])

times_vec = []
times_loop = []

for n in n_values:
    times_vec.append(time_call(f_vec, n, repeat=3))
    times_loop.append(time_call(f_loop, n, repeat=3))

times_vec = np.array(times_vec)
times_loop = np.array(times_loop)

plt.figure()
plt.plot(n_values, times_vec, marker='o', label='vectorized')
plt.plot(n_values, times_loop, marker='o', label='for-loop')
plt.xlabel('n')
plt.ylabel('time (s)')
plt.title('Runtime comparison')
plt.grid(True)
plt.legend()
plt.show()
\end{lstlisting}

\end{document}